\documentclass[11pt]{amsart}
\usepackage[margin=1.1in]{geometry}
\usepackage{amsmath,amssymb,amsthm}
\usepackage{booktabs}
\usepackage{xcolor}
\usepackage[colorlinks=true,linkcolor=blue!50!black,citecolor=blue!50!black]{hyperref}

\newtheorem{theorem}{Theorem}[section]
\newtheorem{lemma}[theorem]{Lemma}
\newtheorem{proposition}[theorem]{Proposition}
\newtheorem{corollary}[theorem]{Corollary}
\theoremstyle{definition}
\newtheorem{definition}[theorem]{Definition}
\theoremstyle{remark}
\newtheorem{remark}[theorem]{Remark}
\newtheorem{convention}[theorem]{Convention}

\DeclareMathOperator{\Inn}{Inn}
\DeclareMathOperator{\Aut}{Aut}
\DeclareMathOperator{\Out}{Out}
\DeclareMathOperator{\Inndiag}{Inndiag}
\DeclareMathOperator{\Syl}{Syl}
\DeclareMathOperator{\Fix}{Fix}
\DeclareMathOperator{\PSL}{PSL}
\DeclareMathOperator{\PSU}{PSU}
\DeclareMathOperator{\PSp}{PSp}
\DeclareMathOperator{\Sp}{Sp}
\DeclareMathOperator{\SL}{SL}
\DeclareMathOperator{\SU}{SU}
\DeclareMathOperator{\GL}{GL}
\DeclareMathOperator{\GU}{GU}
\DeclareMathOperator{\Sz}{Sz}
\DeclareMathOperator{\ord}{ord}
\newcommand{\N}{\mathbb{N}}
\newcommand{\Z}{\mathbb{Z}}
\newcommand{\F}{\mathbb{F}}
\newcommand{\Pos}{\mathcal{P}}
\newcommand{\vp}[1]{v_{#1}}

\title[Groups factorized by non-conjugate maximal subgroups]{Finite groups
that are the product of every pair of non-conjugate maximal subgroups are
soluble}

\author{Richie Sater}
\address{Independent Researcher, Cary, North Carolina, USA}
\email{[redacted-email]}
\urladdr{[redacted-orcid]}

\date{July 27, 2026}

\subjclass[2020]{Primary 20D40; Secondary 20E28, 20D05}
\keywords{Factorizations of finite groups, maximal subgroups, soluble
groups, simple groups, Kourovka Notebook Problem 10.34}

\begin{document}

\begin{abstract}
We prove that every finite group which coincides with the product of any
two of its non-conjugate maximal subgroups is soluble, answering Problem
10.34 of the Kourovka Notebook (V.~S.~Monakhov, 1986) in the negative.
The proof combines a reduction to groups with a unique minimal normal
subgroup $S^k$, a divisibility criterion that excludes a simple group $S$
for \emph{all} $k\geq 1$ at once from a single pair of automorphism-stable
conjugacy classes of subgroups of $S$, a uniform construction of such
pairs for every infinite family of finite simple groups (with Zsigmondy
primes as the arithmetic obstruction), and machine-generated certificates,
produced in GAP, covering the sporadic groups and all small cases,
including every non-abelian simple group of order at most
$1.05\times 10^{7}$.
\end{abstract}

\maketitle

\section{Introduction}\label{sec:intro}

Problem 10.34 of the Kourovka Notebook, posed by V.~S.~Monakhov in 1986
\cite{Kourovka}, asks:

\begin{quote}
\emph{Does there exist a finite non-soluble group that coincides with the
product of any two of its non-conjugate maximal subgroups?}
\end{quote}

Soluble groups with this property exist in abundance: any group
whose maximal subgroups are all conjugate satisfies it vacuously,
and $S_4 = AB$ for every pair of non-conjugate maximals $A,B$. The content
of the problem is whether the property forces solubility. We prove that it
does.

\begin{theorem}[Main Theorem]\label{thm:main}
Let $G$ be a finite group such that $G=AB$ for every pair of maximal
subgroups $A,B\leq G$ that are not conjugate in $G$. Then $G$ is soluble.
\end{theorem}

Throughout, we say $G$ has \emph{property $\mathrm{P}$} if $G=AB$ for all
non-conjugate maximal $A,B\leq G$.

The almost simple case of Theorem~\ref{thm:main} was proved by Tikhonenko
and Tyutyanov \cite{TT2010}. Our argument is independent of \cite{TT2010}:
the criterion developed in \S\ref{sec:criterion} covers the almost simple
case uniformly with the general one (Theorem~\ref{thm:Dprime}), so the
present paper re-proves that result along the way.

\subsection*{Strategy}
A minimal counterexample $G$ has a unique minimal normal subgroup
$N=S^k$ with $S$ non-abelian simple, and embeds in $\Aut(S)\wr S_k$
(\S\ref{sec:reduction}). For an automorphism-stable conjugacy class $[V]$
of suitable subgroups $V<S$, the normalizer $B_V=N_G(V^k)$ is a maximal
subgroup of $G$ of order exactly $|G/N|\cdot|V|^k$
(\S\ref{sec:machinery}). Property $\mathrm{P}$ forces $G=B_VB_W$ for
distinct such classes, and the resulting integrality condition
$|B_V\cap B_W|\in\Z$ produces, for a well-chosen prime $p$, the
divisibility $p^{dk}\mid |G/N|$ with $d$ independent of $k$, which is
incompatible with $|G/N|$ dividing $x^k\,k!$, where $x=|X/\Inn(S)|$ for
the relevant $X\leq\Aut(S)$ (\S\ref{sec:criterion}). The upshot
(Theorem~\ref{thm:D}) is that a \emph{single} pair of stable classes in
$S$, satisfying a valuation inequality, rules out the socle $S^k$ for
every $k\geq 1$ and every admissible embedding at once.

The problem thus reduces to exhibiting such a pair for every finite simple
group $S$. By the classification of finite simple groups this is a finite
list of infinite families plus $26$ sporadic groups and the Tits group.
For the infinite families the pairs are constructed uniformly in
\S\ref{sec:families}, with Zsigmondy primes of the relevant torus
exponents as the obstruction primes; the finitely many exceptional
configurations (Zsigmondy exceptions, small-field maximality exceptions,
the outer automorphisms of $A_6$, and so on) all have order at most
$1.05\times 10^{7}$ apart from two, $\PSL(6,2)$ and $\Omega^+(8,2)$,
which are handled in-text with substitute primes. Everything of order at most
$1.05\times10^{7}$, and every sporadic group, is settled by machine
certificates (\S\ref{sec:machine}).

Three structural facts keep the family analysis manageable. First, when
the two exhibited classes consist of \emph{maximal} subgroups of $S$, all
hypotheses of the machinery are automatic and only stability and the
valuation inequality need checking. Second, no automorphism of a twisted
group of Lie type permutes the types of its twisted building nontrivially
(the twisted groups have no analogue of the untwisted graph
automorphisms, by the structure of their automorphism groups), so their
parabolic classes are stable under every automorphism. The twisted families,
including the triality groups ${}^3D_4(q)$ and the Ree groups, are
consequently among the easiest cases.
Third, when a graph automorphism does fuse the relevant parabolic classes,
the corresponding \emph{flag} parabolics (intersections of fused maximal
parabolics) become poset-maximal among stable classes, and a structural
lemma (Lemma~\ref{lem:P}) shows they satisfy all needed hypotheses; this
is the wreath-product analogue of the ``novelty'' phenomenon for maximal
subgroups of almost simple groups.

\subsection*{Computer verification}
The machine-certified ingredients were produced with GAP~4.16.0
\cite{GAP}, using its simple-group iterators, the character table library
\cite{CTblLib} and \textsc{Atlas} data \cite{ATLAS}. Every certificate is
a single logged line naming the two class orders, the prime, and the
valuations, so each is independently re-checkable by hand or machine. The
scripts and logs, together with an arithmetic verification of the family
proofs of \S\ref{sec:families} over large parameter ranges ($7892$
instances), are publicly available in the GitHub repository
\cite{Repo}. \S\ref{sec:machine} specifies exactly what is claimed from
the machine and how it is certified.

\subsection*{Acknowledgements}
This work was carried out in collaboration with Anthropic's Claude
(2026 Claude~5-family models), which assisted with both the mathematics
and the software. All machine claims are backed by committed logs as
described in \S\ref{sec:machine}, and the author takes full
responsibility for the correctness of every statement in this paper.

\section{Reduction to socle analysis}\label{sec:reduction}

\begin{lemma}[Conjugation lemma]\label{lem:conj}
If $G=AB$ for subgroups $A,B$, then $G=A^xB^y$ for all $x,y\in G$.
\end{lemma}

\begin{proof}
Since $G=AB$, write $xy^{-1}=ab$ with $a\in A$, $b\in B$. Then, as
subsets of $G$,
\[
A^xB^y=x^{-1}Ax\,y^{-1}By=x^{-1}A(ab)By
 =x^{-1}(AB)\,y=x^{-1}Gy=G. \qedhere
\]
\end{proof}

Consequently, property $\mathrm{P}$ depends only on the conjugacy classes
of maximal subgroups: it says every pair of \emph{distinct classes} has
one (hence every) representative pair factorizing $G$.

\begin{lemma}\label{lem:quot}
Property $\mathrm{P}$ is inherited by quotients.
\end{lemma}

\begin{proof}
Maximal subgroups of $G/K$ are the images of the maximal subgroups of $G$
containing $K$; non-conjugate ones lift to non-conjugate ones, and
factorizations project.
\end{proof}

\begin{proposition}[Minimal counterexample structure]\label{prop:min}
Let $G$ be a non-soluble group with property $\mathrm{P}$ of least order.
Then $G$ has a unique minimal normal subgroup $N=S^k$ with $S$ non-abelian
simple and $k\geq 1$, $C_G(N)=1$, $G/N$ is soluble, and $G$ embeds in
$\Aut(S)\wr S_k$ with $N$ mapping onto $\Inn(S)^k$ and $G$ acting
transitively on the $k$ coordinates.
\end{proposition}

\begin{proof}
Every proper quotient of $G$ has property $\mathrm{P}$
(Lemma~\ref{lem:quot}), hence is soluble by minimality of $|G|$. The
soluble radical $R$ of $G$ is trivial: otherwise $G/R$ is non-soluble
(as $R$ is soluble) with property $\mathrm{P}$, contradicting minimality.
So the socle of $G$ is a direct product of non-abelian minimal normal
subgroups. If $N_1\neq N_2$ were two of them, then $G/N_1$ is soluble
while $N_2\cong N_2N_1/N_1\leq G/N_1$ is non-soluble, a contradiction.
Hence there is a unique minimal normal subgroup $N=S^k$, $S$ non-abelian
simple. $C_G(N)$ is a normal subgroup meeting $N$ trivially, so
$C_G(N)=1$, and the conjugation action embeds $G$ into $\Aut(N)=
\Aut(S)\wr S_k$ in the standard way; minimal normality of $N$ forces
transitivity on coordinates. Finally $G/N$ is soluble because it is a
proper quotient.
\end{proof}

\begin{convention}\label{conv:X}
For $G$ as in Proposition~\ref{prop:min} let $G_1$ be the stabilizer of
the first coordinate and $\pi_1\colon G_1\to\Aut(S)$ the projection. Set
$X:=\pi_1(G_1)\cdot\Inn(S)$, the \emph{coordinate closure}; by
coordinate-transitivity all coordinates give the same $X$ up to conjugacy,
and after conjugating inside $\Aut(S)\wr S_k$ we may and do assume
$G\leq X\wr S_k$ with $\Inn(S)\leq X\leq\Aut(S)$. (For the
normalization: fix $g_j\in G$ carrying coordinate $1$ to coordinate $j$,
with $g_1=1$, and let $a_j\in\Aut(S)$ be its $1\!\to\!j$ component. If
$g\in G$ carries coordinate $i$ to coordinate $j$ through the component
$c\in\Aut(S)$, then $g_i\,g\,g_j^{-1}$ stabilizes coordinate $1$ with
component $a_i\,c\,a_j^{-1}$, which therefore lies in
$\pi_1(G_1)\leq X$; so conjugating $G$ by the diagonal-type element
$(a_1,\dots,a_k)\in\Aut(S)^k$ places every component of every element
of $G$ in $X$, and $a_1=1$ leaves $X$ itself unchanged.) Write
\[
x:=|X/\Inn(S)|,\qquad t:=|G/N| .
\]
Then $t$ divides $x^k\,k!$. Since property $\mathrm{P}$ and all criteria
below are invariant under replacing $G$ by a conjugate under
$(a,\dots,a)\in\Aut(S)^k$, statements may be proved after normalizing $X$
up to $\Aut(S)$-conjugacy.
\end{convention}

The following observation is not needed for the proof of
Theorem~\ref{thm:main} but explains why the pairs used later are of
``supplement'' type; it also matches the exhaustive machine data of
\S\ref{sec:machine}.

\begin{remark}[Top--supplement pairs always factorize]
If $M\supseteq N$ is maximal and $B$ is maximal with $B\not\supseteq N$,
then $BN=G$ ($BN$ is a subgroup strictly containing the maximal $B$),
and since $N\leq M$ we get $MB\supseteq NB=G$. Hence property
$\mathrm{P}$ for $G$ is equivalent to: $\mathrm{P}(G/N)$ together with the
factorization of every pair of non-conjugate maximal subgroups not
containing $N$. In roughly $200$ exhaustively tested candidate groups the
machine found no failing top--supplement pair, consistent with this.
\end{remark}

\section{Product supplements}\label{sec:machinery}

Fix $G$ as in Convention~\ref{conv:X}: $N=S^k\leq G\leq X\wr S_k$,
$t=|G/N|$, $x=|X/\Inn(S)|$. For $V\leq S$ write $[V]$ for its
$S$-conjugacy class. An $S$-class is \emph{$X$-stable} if $X$ permutes it,
i.e.\ $V^a$ is $S$-conjugate to $V$ for every $a\in X$ ($\Inn$-stability
being automatic).

\begin{definition}\label{def:poset}
$\Pos_X(S):=\{\,[V] : 1<V<S,\ N_S(V)=V,\ [V]\ X\text{-stable}\,\}$,
partially ordered by containment up to $S$-conjugacy
($[V]\leq[U]$ iff $V\leq U^s$ for some $s\in S$). Call $V$
\emph{normally saturating} if $\langle V^W\rangle=W$ for every overgroup
$V\leq W\leq S$.
\end{definition}

\begin{remark}\label{rem:maxauto}
If $V$ is a \emph{maximal} subgroup of $S$ then, whenever $[V]$ is
$X$-stable: $[V]\in\Pos_X(S)$ (maximal subgroups of a non-abelian simple
group are self-normalizing), $[V]$ is a maximal element of the poset, and
$V$ is normally saturating (its only overgroups are $V$ and $S$, and
$\langle V^S\rangle=S$ by simplicity). So for maximal classes the lemmas
below apply with no further hypotheses.
\end{remark}

\begin{lemma}[Existence]\label{lem:A}
Let $[V]\in\Pos_X(S)$ and set $B_V:=N_G(V^k)$. Then
\begin{enumerate}
\item $B_V\cap N=N_S(V)^k=V^k$;
\item $B_VN=G$;
\item $|B_V|=t\cdot|V|^k$ exactly.
\end{enumerate}
\end{lemma}

\begin{proof}
(i) An element of $N=S^k$ normalizes $V^k$ iff each coordinate normalizes
$V$. (ii) Frattini-type argument: for $g\in G$ write $g=(a_1,\dots,a_k)
\sigma$ with $a_i\in X$, $\sigma\in S_k$. Then $(V^k)^g=V^{a_{1'}}\times
\cdots\times V^{a_{k'}}$ (indices permuted by $\sigma$), and $X$-stability
gives $s_i\in S$ with $V^{a_i}=V^{s_i}$. So $(V^k)^g$ is $N$-conjugate to
$V^k$: there is $n\in N$ with $(V^k)^{gn}=V^k$, i.e.\ $gn\in B_V$, so
$g\in B_VN$. (iii) $|B_V|=|B_VN|\cdot|B_V\cap N|/|N|=t\,|V|^k$.
\end{proof}

\begin{lemma}[Maximality]\label{lem:B}
If in addition $[V]$ is a maximal element of $\Pos_X(S)$ and $V$ is
normally saturating, then $B_V$ is maximal in $G$.
\end{lemma}

\begin{proof}
Let $B_V\leq H\leq G$. Since $H\supseteq B_V$ and $B_VN=G$ we have
$HN=G$, so $H$ surjects onto $G/N$ and in particular acts
coordinate-transitively.

\emph{Step 1 ($H$'s twists recover $X$).} Let $H_1,G_1$ be the
coordinate-$1$ stabilizers. For $g\in G_1$ write $g=hn$ ($h\in H$,
$n\in N$); $n$ fixes every coordinate, hence so does $h$, so $h\in H_1$
and $\pi_1(G_1)=\pi_1(H_1)\cdot\pi_1(N\cap G_1)=\pi_1(H_1)\cdot\Inn(S)$.
Therefore $\pi_1(H_1)$ and $\Inn(S)$ together generate $X$, and a class
of subgroups of $S$ is $X$-stable iff it is stable under $\pi_1(H_1)$.

\emph{Step 2 (Goursat + saturation: $H\cap N$ is a product).} Let
$T:=H\cap N\supseteq V^k$, let $W_i:=\pi_i(T)\leq S$ be the coordinate
projections and $K_i:=T\cap S_i$ the coordinate kernels ($S_i$ the $i$-th
factor). Then $V\leq K_i$ (as $V^k\cap S_i=V$ in coordinate $i$),
$K_i\trianglelefteq W_i$, and $V\leq W_i$. Since $K_i$ is a normal
subgroup of $W_i$ containing $V$, it contains
$\langle V^{W_i}\rangle=W_i$ by saturation. So $K_i=W_i$ for every $i$
and $T=\prod_i W_i$.

\emph{Step 3 (the $W_i$ force $H=B_V$ or $H=G$).} $T$ is $H$-invariant;
$H$ permutes the coordinates transitively and twists by $\pi_1(H_1)$, so
all $[W_i]$ coincide with a single class $[W]$ which, by Step~1, is
$X$-stable. If $W=S$ then $T=N$, so $H\supseteq N$ and $H=HN=G$.
Otherwise $1<V\leq W<S$. Form the normalizer tower
$W\leq N_S(W)\leq N_S(N_S(W))\leq\cdots$; it is strictly increasing
until it becomes self-normalizing, and it cannot reach $S$ (the
penultimate term would be a proper normal subgroup of the simple group
$S$). Its terminal member $U$ is proper, self-normalizing, and $[U]$ is
$X$-stable ($N_S(W^a)=N_S(W)^a$, so stability propagates up the tower),
i.e.\ $[U]\in\Pos_X(S)$ with $[U]\geq[W]\geq[V]$. Maximality of $[V]$ in
$\Pos_X(S)$ forces $[U]=[V]$; then $|U|=|V|$ and $V\leq W\leq U$ gives
$V=W=U$. So $T=V^k$, hence $|H|=|HN|\cdot|H\cap N|/|N|=t\,|V|^k=|B_V|$
and $H=B_V$.
\end{proof}

\begin{lemma}[Non-conjugacy]\label{lem:C}
Distinct $[V]\neq[W]$ in $\Pos_X(S)$ give non-conjugate $B_V,B_W$.
\end{lemma}

\begin{proof}
If $B_V^g=B_W$ then $(B_V\cap N)^g=B_W\cap N$ (as $N\trianglelefteq G$),
i.e.\ $(V^k)^g=W^k$. Reading off any coordinate, $W$ is the image of $V$
under an element of $X$ composed with a coordinate permutation, so
$[W]=[V^a]=[V]$ for some $a\in X$ by $X$-stability, a contradiction.
\end{proof}

The next lemma supplies, structurally, all hypotheses of
Lemma~\ref{lem:B} for the non-maximal (``novelty'') classes used in
\S\ref{sec:families}: the standard parabolic subgroups whose maximal
parabolic overgroups are fused by a graph automorphism. We use the
notation of \cite{Carter}: $S$ an untwisted simple group of Lie type
with split $BN$-pair of rank $\ell\geq2$ (all our applications are
untwisted), $B$ a Borel subgroup,
$\Delta$ the node set of the Dynkin diagram, and for
$J\subseteq\Delta$, $Q_J\supseteq B$ the standard parabolic whose Levi
subgroup $L_J$ is generated by the maximal torus $T$ and the root
subgroups of $J$; the maximal parabolic omitting node $i$ is
$P_i=Q_{\Delta\setminus\{i\}}$.

\begin{lemma}[Parabolic novelties]\label{lem:P}
Fix $J\subseteq\Delta$ and put $V:=Q_J$. Suppose $X$ is such that
\begin{enumerate}
\item[(a)] the $S$-class $[Q_J]$ is $X$-stable, and
\item[(b)] for every $J\subsetneq K\subsetneq\Delta$ the class $[Q_K]$
is \emph{not} $X$-stable.
\end{enumerate}
Then $[V]\in\Pos_X(S)$, $[V]$ is a maximal element of $\Pos_X(S)$, and
$V$ is normally saturating. Consequently Lemmas~\ref{lem:A}--\ref{lem:C}
apply to $V$ with no computational hypothesis.
\end{lemma}

\begin{proof}
Parabolic subgroups are self-normalizing \cite[Thm.~8.4.4]{Carter}, so
$[V]\in\Pos_X(S)$ by (a).

\emph{Overgroups.} Every subgroup containing $V$ contains $B$, and every
subgroup containing a Borel subgroup is a parabolic $Q_K$, $K\supseteq J$,
of the same $BN$-pair \cite[Thms.~8.3.2, 8.4.3]{Carter}; moreover $Q_K$
and $Q_{K'}$ are $S$-conjugate iff $K=K'$ (the type of a parabolic is a
conjugation invariant of the building).

\emph{Poset-maximality.} If $[U]\in\Pos_X(S)$ with $V\leq U^s<S$ for some
$s$, then $U^s=Q_K$ for some $J\subseteq K\subsetneq\Delta$; $K=J$ gives
$[U]=[V]$, while $J\subsetneq K$ contradicts (b), since $[U]=[Q_K]$ would
be $X$-stable.

\emph{Saturation.} Let $V\leq W\leq S$, so $W=Q_K$ ($K\supseteq J$) or
$W=S$; we must show $\langle V^W\rangle=W$. The normal closure of $V$ in
$W$ contains that of $B$. Write $W=U_K\rtimes L_K$ (Levi decomposition;
for $W=S$ read $U_K=1$, $L_K=S$). Then $U_K\leq B$, and modulo $U_K$ the
image of $B$ is a Borel subgroup $B_L=T\cdot(U\cap L_K)$ of $L_K$
containing the full torus $T$. For each node $\alpha\in K$ the
reflection representative $n_\alpha$ lies in $L_K$, and
$U_{-\alpha}=U_\alpha^{\,n_\alpha}$. The subgroup
$\langle B_L^{\,L_K}\rangle$ contains every $L_K$-conjugate of every
subgroup of $B_L$; in particular it contains $T$, each $U_\alpha$
($\alpha\in K$ positive), and each $U_{-\alpha}=U_\alpha^{\,n_\alpha}$.
These generate $L_K$ \cite[Thm.~8.1.5 with \S7.2]{Carter}. So
$\langle B^W\rangle\supseteq U_K\cdot L_K=W$.
\end{proof}

\section{The divisibility criterion}\label{sec:criterion}

For a prime $p$, $\vp{p}$ denotes the $p$-adic valuation and $s_p(k)$ the
digit sum of $k$ in base $p$, so that $\vp{p}(k!)=(k-s_p(k))/(p-1)$.

\begin{theorem}[Divisibility criterion]\label{thm:D}
Let $S,X$ be as in Convention~\ref{conv:X} and let $[V]\neq[W]$ be
maximal elements of $\Pos_X(S)$, both normally saturating. If some prime
$p$ satisfies
\[
d:=\vp{p}(|S|)-\vp{p}(|V|)-\vp{p}(|W|)\;>\;\vp{p}(x),
\]
then no group $G$ with socle $S^k$ and coordinate closure $X$, for any
$k\geq2$, has property $\mathrm{P}$.
\end{theorem}

\begin{proof}
By Lemmas~\ref{lem:A}--\ref{lem:C}, $B_V$ and $B_W$ are non-conjugate
maximal subgroups of $G$, so property $\mathrm{P}$ forces $G=B_VB_W$ and
\[
|B_V\cap B_W|=\frac{|B_V|\,|B_W|}{|G|}
 = t\cdot\Bigl(\frac{|V|\,|W|}{|S|}\Bigr)^{k}\in\Z ,
\]
whence $p^{dk}\mid t$. But
\[
\vp{p}(t)\;\leq\;k\,\vp{p}(x)+\vp{p}(k!)
 \;=\;k\,\vp{p}(x)+\frac{k-s_p(k)}{p-1}
 \;<\;k\bigl(\vp{p}(x)+1\bigr)\;\leq\;dk ,
\]
a contradiction.
\end{proof}

\begin{remark}
(i) The criterion is independent of $k$: this is what closes infinite
families, which no counting bound can do, since $(k!)^{1/k}\to\infty$
outgrows any fixed ratio $|V||W|/|S|$. (ii) Only the \emph{existence} of
$B_V,B_W$ is used; no classification of the maximal subgroups of $G$
(or even of $S$) is needed. In particular the criterion is insensitive
to completeness of tabulated maximal-subgroup data: it needs only that
the two exhibited subgroups are maximal (or Lemma~\ref{lem:P}-novelties)
and their classes stable.
\end{remark}

\begin{theorem}[$k=1$: the almost simple case]\label{thm:Dprime}
The same criterion excludes the almost simple case: for
$S\leq G\leq\Aut(S)$ put $X:=G$, $x:=|G/S|=t$. For $[V]$ a maximal
element of $\Pos_X(S)$, normally saturating, $B_V:=N_G(V)$ satisfies
$B_V\cap S=V$, $B_VS=G$, $|B_V|=t|V|$, and $B_V$ is maximal in $G$.
Hence if $[V]\neq[W]$ are two such classes and
$d=\vp{p}(|S|)-\vp{p}(|V|)-\vp{p}(|W|)>\vp{p}(x)$ for some prime $p$,
then $G$ does not have property $\mathrm{P}$.
\end{theorem}

\begin{proof}
$B_VS=G$ is the Frattini argument via $X$-stability as in
Lemma~\ref{lem:A}; $B_V\cap S=N_S(V)=V$; the maximality of $B_V$ is the
proof of Lemma~\ref{lem:B} without the Goursat step: $T:=H\cap S$ is an
$H$-invariant subgroup containing $V$; if $T=S$ then $H\supseteq S$ and
$H=HS=G$; otherwise the normalizer tower of $T$ ends at a proper
self-normalizing subgroup whose class is $X$-stable (here
$G=B_VS\leq HS$ forces $HS=G$, so $X=G$ is generated by the images of
$H$ and $S$, and stability may be tested against the $H$-invariance of
$T$, exactly as in Step~1 of Lemma~\ref{lem:B}) and lies above $[V]$,
so poset-maximality forces $T=V$ and $H\leq N_G(V)=B_V$. Then
$G=B_VB_W$ forces $t\,|V||W|/|S|\in\Z$, i.e.\ $p^d\mid t=x$,
contradicting $d>\vp{p}(x)$.
\end{proof}

\begin{convention}\label{conv:excluded}
From now on, ``$S$ is \emph{excluded}'' means: for every $X$ with
$\Inn(S)\leq X\leq\Aut(S)$ there exist classes as in
Theorem~\ref{thm:D}/\ref{thm:Dprime}; consequently no finite group $G$
with unique minimal normal subgroup $S^k$ ($k\geq1$, any admissible
embedding) has property $\mathrm{P}$, i.e.\ $S$ cannot occur in a
minimal counterexample to Theorem~\ref{thm:main}.
\end{convention}

\section{The machine base}\label{sec:machine}

\begin{proposition}\label{prop:base}
Every non-abelian simple group $S$ with $|S|\leq 1.05\times10^{7}$ is
excluded (in the sense of Convention~\ref{conv:excluded}).
\end{proposition}

The proof of Proposition~\ref{prop:base} is a finite computation in
GAP~4.16.0 \cite{GAP}, organized as follows; every claim is backed by a
committed log file, and the group inventory is cross-checked against
GAP's \texttt{SimpleGroupsIterator} in two installments, so that no
group in range is skipped: the $47$ groups with $|S|<5\times10^5$
(receipt \texttt{verify\_coverage.log}) and the $51$ groups with
$5\times10^5\leq|S|\leq1.05\times10^7$ (receipt
\texttt{verify\_coverage\_big.log}, which also re-checks that the only
two survivors of the large-range maximal-pair sweep, $\Sp(4,4)$ at
$x=4$ and $\PSL(5,2)$ at $x=2$, are closed by the novelty certificates
of item~(2)).

\begin{enumerate}
\item \emph{Maximal-class pairs (sweeps J, J2, J4, J3, J5, J6).} For each
simple $S$ in range and each $X$ with $\Inn(S)\leq X\leq\Aut(S)$, the
program lists the $X$-stable classes of maximal subgroups of $S$, and
searches pairs $([V],[W])$ and primes $p$ satisfying the inequality of
Theorem~\ref{thm:D}. For all but six socles in range this succeeds for
every $X$. A logged certificate line records
$(|V|,|W|,p,d)$; e.g.\ for $S=A_5$: $(|A_4|,|S_3|,5,1)$, excluding
$A_5^k$ for all $k$ and all $X$, the first unbounded-$k$ closure. In
addition, the alternating groups $A_{11}$--$A_{14}$, whose orders
exceed the base bound, carry certificates of the same form for every
$X$ (sweep J5, log \texttt{sweepJ5\_smallAn.log}); Theorem~\ref{thm:an}
uses these for $11\leq n\leq14$.
\item \emph{Novelty pairs (sweeps K, K2, K3, K4).} Six $(S,X)$ families
admit too few stable maximal classes, all for the expected reason (an
outer automorphism fuses classes): $\PSL(3,2)$ and $\PSL(2,11)$ at
$x=2$; $A_6$ for $X/\Inn\in\{2_2,2_3,2^2\}$; $\PSL(3,4)$ for several
$X$; $\Sp(4,4)$ at $x=4$; $\PSL(5,2)$ at $x=2$. Each is closed by a pair
involving non-maximal classes that are maximal in $\Pos_X(S)$, with
normal saturation verified computationally (sweep K2) or structurally
(Lemma~\ref{lem:P}). The K2 saturation certificates operate at the
level of conjugacy classes, not just subgroup orders: the log records,
per group, every self-normalizing class (listed by order) and every
saturation failure. The recorded failures (order-$6$ classes of
$A_6$ and $\PSL(3,4)$, and order-$12$ classes of $\PSL(3,4)$) belong
to classes used in no exclusion pair; every class actually used either
lies in a group all of whose self-normalizing classes saturate
($\PSL(3,2)$, and $\PSL(2,11)$, where this matters because its pair
$(A_4,D_{12})$ has two distinct classes of the same order $12$) or has
order distinct from every failing class ($A_6$: orders $8,36$;
$\PSL(3,4)$: orders $72,192$). Sample certificates: $\PSL(3,2)$ at
$x=2$ via $(S_3,\,7{:}3)$ at $p=2$, $d=2>1$; $\Sp(4,4)$ at $x=4$ via
(Borel,
subfield $\Sp(4,2)$) at $p=5$ and $17$; $\PSL(5,2)$ at $x=2$ via the
incident-flag parabolics $(Q_{\{2,3\}},Q_{\{1,4\}})$ (the
point--hyperplane and line--$3$-space flag stabilizers, in the standard
parabolic notation of Lemma~\ref{lem:P}) at $p=5$ and $31$. The last
two are
instances of the uniform constructions of Theorem~\ref{thm:graph} below.
\item \emph{Consistency checks.} Independently of the criterion, property
$\mathrm{P}$ was tested exhaustively (all admissible $G$, transitive
tops) for: all socles $S^2$ with $|S|<5\times10^5$; $S^3$ for the
smallest socles and for the eleven socles in that range surviving a
$k=3$ prefilter (complete except over $\PSU(3,5)^3$, where $20$ of $34$
candidate tops have been checked at the time of writing);
$S^4$ for small $S$; all non-soluble groups of order $\leq2000$;
all perfect groups of order $\leq2\times10^6$. No group with property
$\mathrm{P}$ and non-abelian socle was found, and every observed failing
pair was of supplement--supplement or top--top type, consistent with
the theory above.
\end{enumerate}

\begin{proposition}\label{prop:sporadic}
Every sporadic simple group, and the Tits group ${}^2F_4(2)'$, is
excluded.
\end{proposition}

\begin{proof}
Six of the $26$ sporadic groups ($M_{11}$, $M_{12}$, $M_{22}$,
$M_{23}$, $J_1$, $J_2$) have order at most $1.05\times10^7$ and fall
under Proposition~\ref{prop:base}. The remaining $20$ ($M_{24}$, $J_3$,
$J_4$, $HS$, $McL$, $Co_1$, $Co_2$, $Co_3$, $Suz$, $He$, $Ru$, $O'N$,
$\mathrm{Fi}_{22}$, $\mathrm{Fi}_{23}$, $\mathrm{Fi}_{24}'$, $HN$,
$Ly$, $Th$, $B$, $M$) and the Tits group are handled by a dedicated
computation (sweep~M): stable maximal classes and pair searches use the
\textsc{Atlas} maximal-subgroup data via the GAP character table
library \cite{CTblLib,ATLAS}. Since $|\Out(S)|\leq2$ for all these
groups, stability is either automatic ($\Out=1$, or a unique class of
its order) or checked against the stored class fusions. Sample
certificates: the Monster via $(59{:}29,\,41{:}40)$ at $p=71$; the
Baby Monster at $p=31$; $\mathrm{Fi}_{24}'$ at $p=23$. By the remark
after Theorem~\ref{thm:D}, the criterion does not require completeness
of the \textsc{Atlas} maximal-subgroup lists: only that the two
subgroups used are maximal and their classes stable, which for these
tables is established data.
\end{proof}

Beyond the base, the family proofs of \S\ref{sec:families} carry their
own numeric receipts: the valuation inequalities of
Theorems~\ref{thm:psl2}--\ref{thm:graph} were verified mechanically over
large parameter ranges ($1272$ prime powers for Theorem~\ref{thm:psl2};
$n\leq10^4$ for Theorem~\ref{thm:an}; $7892$ parameter instances across
the Lie families, ranks up to $25$ and fields up to $3000$, for
Theorems~\ref{thm:twisted2}--\ref{thm:graph}), with zero failures and
with the Zsigmondy-exception set exactly as documented in the proofs.

\section{The infinite families}\label{sec:families}

Throughout this section $S$ is simple, $X$ is any group with
$\Inn(S)\leq X\leq\Aut(S)$, and $x=|X/\Inn(S)|$. By
Remark~\ref{rem:maxauto}, when the two exhibited classes consist of
maximal subgroups of $S$, only their $X$-stability and the valuation
inequality require proof. We use Zsigmondy's theorem \cite{Zsigmondy}
throughout: for integers $a\geq2$, $e\geq3$ with $(a,e)\neq(2,6)$, there
is a prime $r$ with $\ord_r(a)=e$ (a \emph{primitive prime divisor} of
$a^e-1$); such $r$ satisfies $r\equiv1\pmod e$. For $q=p^f$ we write
$\Phi_e=\Phi_e(q)$ for cyclotomic values and always apply Zsigmondy to
the prime $p$ with exponent $ef$, so $\ord_r(p)=ef$ and
$r\equiv1\pmod{ef}$.

\subsection{Projective special linear groups of rank one}

\begin{theorem}\label{thm:psl2}
$S=\PSL(2,q)$, $q=p^f\geq4$, is excluded.
\end{theorem}

\begin{proof}
For $q\leq271$, $q\neq256$ (in particular every exceptional case of
Dickson's subgroup theorem) this is machine-certified
(Proposition~\ref{prop:base}); $q=256$ exceeds the machine base
($|\PSL(2,256)|=16776960>1.05\times10^7$) and is covered by Case~B
below, which is independent of the base. Assume $q\geq13$. Recall
$|S|=q(q^2-1)/e$ with $e=\gcd(2,q-1)$, and
$\Out(S)=C_e\times C_f$, so $x\mid ef$: that is, $x\mid2f$ for $q$ odd
and $x\mid f$ for $q$ even.

\emph{Case A: $q$ odd, $q\geq13$.} Let $U:=N_S(T_+)$ and $V:=N_S(T_-)$
be the normalizers of the split and nonsplit maximal tori: dihedral of
orders $q-1$ and $q+1$, and maximal in $S$ for $q\geq13$ by Dickson's
theorem \cite{Dickson,Huppert}. All split (resp.\ nonsplit) maximal tori
are $S$-conjugate and every automorphism maps one to one (there is no
graph automorphism in type $A_1$; diagonal and field automorphisms act
algebraically), so $[U]$ and $[V]$ are $X$-stable for every $X$. Now
$|U||V|/|S|=e/q$, so for the defining prime $p$:
$d=\vp{p}(|S|)-0-0=f$, while $x\mid2f$ gives
$\vp{p}(x)\leq\vp{p}(f)\leq\log_pf<f=d$. Theorem~\ref{thm:D}/%
\ref{thm:Dprime} applies.

\emph{Case B: $q=2^f$, $f\geq3$.} Here $|S|=q(q^2-1)$ and $x\mid f$.
Let $U$ be a Borel subgroup (the normalizer of a Sylow $2$-subgroup), of
order $q(q-1)$, and $V:=N_S(T_+)$, dihedral of order $2(q-1)$, maximal
for $q\geq8$ \cite{Dickson,Huppert}. $[U]$ is the class of Sylow-$2$
normalizers and $[V]$ the unique class of split torus normalizers: both
$X$-stable. Now $|U||V|/|S|=2(q-1)/(q+1)$ with $q+1$ odd and coprime to
$2(q-1)$, so it suffices to find a prime $r\mid q+1$ with
$\vp{r}(q+1)>\vp{r}(f)$. If $f\neq3$, take $r$ a primitive prime divisor
of $2^{2f}-1$: then $r\mid(2^{2f}-1)/(2^f-1)=q+1$ and
$r\equiv1\pmod{2f}$, so $r>f$ and $\vp{r}(x)=0<\vp{r}(q+1)=d$. If $f=3$
($q=8$): $q+1=9$, $r=3$: $d=2>1\geq\vp{3}(x)$.
\end{proof}

\subsection{Alternating groups}

\begin{theorem}\label{thm:an}
$S=A_n$, $n\geq5$, is excluded.
\end{theorem}

\begin{proof}
For $n\leq14$ this is machine-certified: $n\leq10$ falls under
Proposition~\ref{prop:base} (including $n=6$, whose exceptional outer
automorphisms fuse classes; all four $X$ are handled there), and
$11\leq n\leq14$ by the sweep-J5 certificates of \S\ref{sec:machine}.
Assume $n\geq15$, so $\Aut(A_n)=S_n$ and $x\mid2$.

Let $p$ be the largest prime $\leq n$. For $m<n/2$ let
$M_m:=(S_m\times S_{n-m})\cap A_n$, and for $n$ even let
$I:=(S_{n/2}\wr S_2)\cap A_n$. Each $M_m$ is maximal in $A_n$, and $I$
is maximal outside a finite list of exceptions all with $n\leq14$
\cite{LPS1987}. Conjugation by $S_n$ preserves the partition shape
(resp.\ block structure), each shape is a single $A_n$-class, and
$\Aut(A_n)=S_n$ for $n\neq6$: all these classes are $X$-stable for every
$X$.

For any two distinct classes among
$\{M_m:\,n-p<m<n/2\}\cup\{I\ \text{if $n$ even and}\ n/2<p\}$:
$\vp{p}(|A_n|)=1$ (Bertrand: $n/2<p\leq n$), while $\vp{p}$ of each
subgroup order is $0$ (all parts of the partitions involved are $<p$),
and $\vp{p}(x)\leq\vp{p}(2)=0$. So $d=1>0$ and
Theorem~\ref{thm:D}/\ref{thm:Dprime} applies, provided two such classes
exist. The number of admissible intransitive $m$ is
$p-\lfloor n/2\rfloor-1$, plus one for $I$ when $n$ is even; two classes
exist whenever $p\geq n/2+3$. By Nagura's theorem \cite{Nagura} (a prime
in $(x,6x/5]$ for $x\geq25$), for $n\geq31$ there is a prime
$p\in(5n/6,n]$, and $5n/6>n/2+3$ for $n>18$. For $15\leq n\leq30$ the
largest prime $\leq n$ satisfies the bound by inspection
($p=13,13,17,17,19,19,19,19,23,23,23,23,23,29,29,29$ for
$n=15,\dots,30$).
\end{proof}

\subsection{Lie type without graph symmetry}

\begin{theorem}\label{thm:nograph}
Let $S$ be a simple group of Lie type of rank $\geq2$ over $\F_q$,
$q=p^f$, of a type with no symmetry of its Dynkin diagram in the given
characteristic:
\[
C_n=\PSp(2n,q)\ (n\geq2,\ p\ \text{odd if}\ n=2),\quad
B_n=\Omega(2n{+}1,q)\ (n\geq3,\ q\ \text{odd}),
\]
\[
G_2(q)\ (p\neq3),\quad F_4(q)\ (p\neq2),\quad E_7(q),\quad E_8(q).
\]
Then $S$ is excluded.
\end{theorem}

\begin{proof}
Write $h$ for the relevant torus exponent: $h=2n$ for $B_n/C_n$, and
$h=6,12,18,30$ for $G_2,F_4,E_7,E_8$.

Let $P_1,P_2$ be maximal parabolic subgroups from two distinct classes
(rank $\geq2$ guarantees these; take the ends of the diagram). They are
maximal subgroups of $S$: any overgroup of a parabolic contains a Borel
subgroup, and every subgroup containing a Borel is parabolic
\cite[Thms.~8.3.2, 8.4.3]{Carter}.

\emph{Stability.} $\Aut(S)$ is generated by inner, diagonal, field and
graph automorphisms \cite{Steinberg}, \cite[Ch.~12]{Carter}. Diagonal
and field automorphisms preserve types of vertices of the building,
hence each parabolic class; type-permuting automorphisms arise only from
diagram symmetries, absent for the listed types in the listed
characteristics (the exceptional symmetries of $B_2/C_2$, $F_4$ in
characteristic~$2$ and $G_2$ in characteristic~$3$ are excluded by
hypothesis and treated in Theorem~\ref{thm:graph}). So every parabolic
class is $X$-stable for every $X$.

\emph{The obstruction prime.} Let $r$ be a primitive prime divisor of
$p^{hf}-1$. Then:
(i) $r\mid|S|$, since the order formula of each listed type contains the
factor $q^h-1$ (equivalently $\Phi_h(q)$) \cite[\S14.1]{Carter}.
(ii) $r$ divides no maximal parabolic order: $|P_i|=q^{N_i}|L_i|
(q-1)^{c_i}$ with $L_i$ the semisimple part of the Levi factor, whose
order is a product of a $q$-power and factors $q^j-1$ with $j\leq h_i$,
$h_i$ the largest torus exponent of $L_i$; in every listed case every
proper Levi has $h_i<h$ ($B_n/C_n$: Levi types
$A_{i-1}\times C_{n-i}$ give $j\leq\max(i,2(n-i))<2n$; $G_2$: Levis
$A_1$, $j\leq2$; $F_4$: Levis $B_3,C_3,A_2{\times}A_1$, $j\leq6$; $E_7$:
Levis $E_6,D_6,A_6,\dots$, exponents $\leq12$; $E_8$: Levis
$E_7,D_7,A_7$, exponents $\leq18$). By primitivity $r$ divides none of
these, nor $q$, nor $q-1$.
(iii) $r\nmid x$: $r\equiv1\pmod{hf}$ with $h\geq4$, so
$r>hf\geq4f\geq d_S f\geq x$, where $d_S\leq\gcd(2,q-1)^{\,2}\leq4$
bounds the diagonal part and there is no graph part.

So $d=\vp{r}(|S|)\geq1>0=\vp{r}(x)$ for the pair $([P_1],[P_2])$, and
Theorem~\ref{thm:D}/\ref{thm:Dprime} applies, provided $r$ exists.

\emph{Zsigmondy exceptions.} $p^{hf}=2^6$ forces $p=2$ and $hf=6$:
either $G_2(2)$ (not simple; $G_2(2)'\cong\PSU(3,3)$ is in the machine
base) or $\Sp(6,2)$ (machine base). For $h\in\{12,18,30\}$, $hf=6$ is
impossible; $B_n$ has $q$ odd.
\end{proof}

\subsection{Twisted groups of twisted rank at least two}

\begin{theorem}\label{thm:twisted2}
Let $S$ be one of
\[
\PSU(n,q)\ (n\geq4),\quad
P\Omega^-(2n,q)={}^2D_n(q)\ (n\geq4),\quad
{}^3D_4(q),\quad {}^2E_6(q),\quad
{}^2F_4(q)\ (q=2^f,\ f\ \text{odd}\geq3).
\]
Then $S$ is excluded.
\end{theorem}

\begin{proof}
These are exactly the twisted types of twisted $BN$-pair rank $\geq2$,
so $S$ has at least two classes of maximal parabolic subgroups
$P_1,P_2$, maximal in $S$ as in Theorem~\ref{thm:nograph} (the cited
$BN$-pair theorems hold for the twisted groups).

\emph{Stability.} We claim every parabolic class is $X$-stable for every
$X$. Every $a\in\Aut(S)$
permutes the parabolic subgroups: they are exactly the overgroups of the
normalizers of the Sylow $p$-subgroups (Borel subgroups), an
automorphism-invariant description; hence $a$ permutes the finitely many
parabolic classes, preserving inclusion up to conjugacy, i.e.\ acts on
the type set of the twisted building. Now use the structure of
$\Aut(S)$: for the twisted types, $\Aut(S)$ is $\Inndiag(S)$ extended by
the cyclic group generated by a field automorphism, and there is no
analogue of the graph automorphisms of the untwisted groups (nor of the
exceptional graph automorphisms of $B_2$, $G_2$, $F_4$)
\cite[Thm.~2.5.12]{GLS3}. Inner automorphisms fix each class. In their
standard representatives, diagonal and field automorphisms normalize
the standard Borel subgroup and maximal torus and act on the twisted
root system preserving each simple type, so they map each standard
parabolic $Q_J$ to itself. Hence every automorphism fixes every
parabolic class. (For the rank-two group ${}^2F_4(q)$ one can also see
directly that the two maximal parabolic classes cannot be interchanged:
their Levi factors ${}^2B_2(q)$ and $\SL_2(q)$ are non-isomorphic. Note
that a diagram-asymmetry argument alone would not suffice at twisted
rank two: the twisted diagrams of $\PSU(4,q)$ and $\PSU(5,q)$ have a
symmetric underlying graph, and the untwisted example $\Sp(4,2^f)$
shows that a multiple bond by itself does not prevent a type-swapping
automorphism.) So every parabolic
class is $X$-stable for every $X$.

\emph{The obstruction prime.} Choose the exponent $e$ and $r$ a
primitive prime divisor of $p^{e}-1$ per type; in each case $r$ divides
$|S|$ but neither $|P_1|$ nor $|P_2|$:
\begin{itemize}
\item $\PSU(n,q)$, $n$ odd: $e=2nf$, so $r\mid\Phi_{2n}(q)\mid q^n+1$,
which divides $|\SU(n,q)|=q^{n(n-1)/2}\prod_{i=2}^n(q^i-(-1)^i)$. The
Levi factor of $P_i$ ($i=1,2$) is a central product of $\GL_i(q^2)$ and
$\GU(n-2i,q)$: its $p'$-order is a product of factors $q^{2j}-1$
($j\leq2$) and $q^j-(-1)^j$ ($j\leq n-2$), all dividing $p^m-1$ with
$m\leq\max(4f,2(n-2)f)<2nf$.
\item $\PSU(n,q)$, $n$ even: $e=2(n-1)f$,
$r\mid\Phi_{2(n-1)}(q)\mid q^{n-1}+1$ (the $i=n-1$ factor). Levi
exponents are $\leq\max(4f,2(n-3)f,(n-2)f)<2(n-1)f$ for $n\geq4$.
\item ${}^2D_n(q)$, $n\geq4$: $e=2nf$, $r\mid\Phi_{2n}(q)\mid q^n+1$,
which divides $q^{n(n-1)}(q^n+1)\prod_{i=1}^{n-1}(q^{2i}-1)$. Levi of
$P_i$: $\GL_i(q)\times O^-(2(n-i),q)$-type, exponents
$\leq\max(2f,2(n-1)f)<2nf$.
\item ${}^3D_4(q)$: $e=12f$, $r\mid\Phi_{12}(q)=q^4-q^2+1$, which
divides $q^8+q^4+1=(q^4+q^2+1)(q^4-q^2+1)$, hence
$|{}^3D_4(q)|=q^{12}(q^8+q^4+1)(q^6-1)(q^2-1)$. The two Levi factors
are $\SL_2(q^3)\cdot(q-1)$ and $\SL_2(q)\cdot(q^3-1)$ (mod centers):
$p'$-orders divide $(q^6-1)(q-1)$ and $(q^2-1)(q^3-1)$, exponents
$\leq6f<12f$.
\item ${}^2E_6(q)$: $e=18f$, $r\mid\Phi_{18}(q)\mid q^9+1$, which
divides $q^{36}(q^{12}-1)(q^9+1)(q^8-1)(q^6-1)(q^5+1)(q^2-1)$. The four
maximal parabolic classes correspond to the four $\tau$-orbits of nodes
of $E_6$ ($\tau$ the order-two symmetry): $\{2\}$, $\{4\}$, $\{1,6\}$,
$\{3,5\}$; the complementary $\tau$-stable subdiagrams are $A_5$,
$A_2{\times}A_1{\times}A_2$, $D_4$, $A_2{+}A_1{+}A_1$ with induced
twists, so the Levi derived groups are $\SU(6,q)$,
$\SL_2(q){\times}\SL_3(q^2)$, ${}^2D_4(q)$-type, and
$\SL_3(q){\times}\SL_2(q^2)$, with cyclotomic exponents
$\leq10f,\,6f,\,8f,\,4f$ respectively, all $<18f$. Any two classes
serve as $P_1,P_2$.
\item ${}^2F_4(q)$, $f$ odd $\geq3$: $e=12f$,
$r\mid\Phi_{12}(q)\mid q^6+1$, which divides
$|{}^2F_4(q)|=q^{12}(q^6+1)(q^4-1)(q^3+1)(q-1)$. The two Levi derived
groups are ${}^2B_2(q)\cong\Sz(q)$ and $\SL_2(q)$
\cite[\S2.9]{GLS3}: $p'$-orders divide $(q^2+1)(q-1)^2$ and
$(q^2-1)(q-1)$, exponents $\leq4f<12f$. (${}^2F_4(2)$ is not simple;
its derived group, the Tits group, is in the machine base.)
\end{itemize}

\emph{$r\nmid x$.} In every case $x$ divides $d_S\cdot c f$ where the
diagonal order $d_S$ divides $q+1$ (types ${}^2A$, ${}^2E_6$) or $4$
(type ${}^2D$) or is $1$ (${}^3D_4$, ${}^2F_4$), and the field part has
order $cf$ with $c\leq3$. Since $r\equiv1\pmod e$ with $e\geq6f$, $r$
exceeds the field-part order; $r\mid d_S$ would give either $r\mid q+1$,
so $\ord_r(p)\mid2f<e$, or $r\leq4<r$, both impossible. Hence
$\vp{r}(x)=0<1\leq d$ and Theorem~\ref{thm:D}/\ref{thm:Dprime} applies.

\emph{Zsigmondy exceptions.} $r$ fails to exist only if $p^e=2^6$. Here
$e\geq6$ always, and $e=6$ occurs only for $\PSU(4,q)$ with $f=1$; then
$p^6=2^6$ forces $S=\PSU(4,2)$, which is in the machine base. No other
case reaches $p^e=2^6$.
\end{proof}

\subsection{Twisted groups of rank one}

\begin{theorem}\label{thm:twisted1}
Let $S$ be one of
\[
\PSU(3,q)\ (q\geq3),\quad
\Sz(q)={}^2B_2(q)\ (q=2^f,\ f\ \text{odd}\geq3),\quad
{}^2G_2(q)\ (q=3^f,\ f\ \text{odd}\geq3).
\]
Then $S$ is excluded.
\end{theorem}

\begin{proof}
Each group has a single class of parabolic subgroups, namely the Borel
subgroups $B=N_S(Q)$, $Q\in\Syl_p(S)$; $B$ is a \emph{maximal}
subgroup (rank-one $BN$-pair) and its class is $X$-stable for every $X$.
We pair it with a second $X$-stable maximal class avoiding a Zsigmondy
prime.

\emph{$\PSU(3,q)$, $q\geq8$.} $|S|=q^3(q^3+1)(q^2-1)/d$,
$d=\gcd(3,q+1)$; $|B|=q^3(q^2-1)/d$. Let $V$ be the stabilizer of a
nonisotropic point of the unitary geometry, of order $q(q-1)(q+1)^2/d$:
maximal for $q\geq3$, $q\neq5$ \cite[Table~8.5]{BHRD} (class
$\mathcal{C}_1$), one class, preserved by every semilinear
automorphism. Take $r$ a primitive prime divisor of $p^{6f}-1$ (exists:
$p^{6f}=2^6$ would force $q=2$, not simple): $r\mid\Phi_6(q)=q^2-q+1$,
which divides $(q^3+1)/(q+1)$, hence $|S|$; $\ord_r(p)=6f$ rules out
$r$ dividing $q(q^2-1)$ or $q(q-1)(q+1)^2$. Also $x\mid2df$, and
$r\equiv1\pmod{6f}$ gives $r>2f$, while $r\mid d\mid q+1$ would force
$\ord_r(p)\mid2f$. So $d_{\mathrm{val}}:=\vp{r}(|S|)\geq1>0=\vp{r}(x)$.
For $q\in\{3,4,5,7\}$ (in particular the maximality exception $q=5$),
$S$ is in the machine base.

\emph{$\Sz(q)$.} $|S|=q^2(q^2+1)(q-1)$; $|B|=q^2(q-1)$. Let
$V:=D_{2(q-1)}$, the normalizer of a cyclic subgroup of order $q-1$:
maximal by Suzuki's classification \cite{Suzuki1962}, one class, stable
under $\Aut(S)=S\rtimes C_f$. Let $r$ be a primitive prime divisor of
$2^{4f}-1$ (exists: $4f\geq12$): $r\mid\Phi_4(q)=q^2+1$, so $r\mid|S|$,
$r\nmid|B|$, $r\nmid|V|$; $x\mid f$ and $r\equiv1\pmod{4f}>f$. So
$\vp{r}(|S|)\geq1>0=\vp{r}(x)$.

\emph{${}^2G_2(q)$.} $|S|=q^3(q^3+1)(q-1)$; $|B|=q^3(q-1)$. Let
$V:=C_S(\iota)\cong\langle\iota\rangle\times\PSL(2,q)$ for an involution
$\iota$ (one class of involutions, so $[V]$ is $\Aut$-stable),
$|V|=q(q^2-1)$; $V$ is maximal \cite{Kleidman1988}. Let $r$ be a
primitive prime divisor of $3^{6f}-1$ (exists: $6f\geq18$):
$r\mid\Phi_6(q)=q^2-q+1$, which divides $(q^3+1)/(q+1)$, so $r\mid|S|$;
$\ord_r(3)=6f$ rules out $r\mid q^3(q-1)$ and $r\mid q(q^2-1)$;
$x\mid f<r$. So $\vp{r}(|S|)\geq1>0=\vp{r}(x)$.
(${}^2G_2(3)$ is not simple; ${}^2G_2(3)'\cong\PSL(2,8)$ is covered by
Theorem~\ref{thm:psl2}.)
\end{proof}

\subsection{Untwisted types with graph symmetry}

\begin{theorem}\label{thm:graph}
Let $S$ be one of
\[
\PSL(n,q)\ (n\geq3),\quad P\Omega^+(2n,q)=D_n(q)\ (n\geq4),\quad
E_6(q),
\]
\[
\Sp(4,q)\ (q=2^f,\ f\geq2),\quad F_4(q)\ (q=2^f),\quad
G_2(q)\ (q=3^f).
\]
Then $S$ is excluded.
\end{theorem}

\begin{proof}
We keep the standard-parabolic notation $Q_J$ of Lemma~\ref{lem:P}. The
\emph{graph part} of $X$ is the image of $X$ in the quotient of
$\Out(S)$ by its inndiag--field part: a subgroup of
$C_2=\langle\delta\rangle$ for $\PSL(n,q)$ ($n\geq3$), $D_n$ ($n\geq5$),
$E_6$; a subgroup of $S_3$ (triality) for $D_4$; while for $\Sp(4,2^f)$,
$F_4(2^f)$, $G_2(3^f)$, $\Out(S)$ is cyclic of order $2f$ generated by
the exceptional graph-field automorphism $\rho$ with $\rho^2=\varphi_p$
\cite[Thm.~2.5.12]{GLS3}, and ``$X$ has graph part'' means the image of
$X$ in $\Out(S)$ is not contained in $\langle\rho^2\rangle$. An
automorphism acts on the building permuting vertex types by a diagram
symmetry; automorphisms with trivial graph part preserve every type. By
Convention~\ref{conv:X} we may normalize the graph part up to
$\Aut(S)$-conjugacy; for $D_4$ this lets us assume a $C_2$ graph part is
generated by the symmetry fixing nodes $1,2$ and swapping $3,4$.

\emph{Case A: $X$ with trivial graph part.} Every parabolic class is
$X$-stable, and the argument of Theorem~\ref{thm:nograph} applies with
these data:
\begin{itemize}
\item $\PSL(n,q)$: pair $([P_1],[P_2])$, $r$ a primitive prime divisor
of $p^{nf}-1$, so $r\mid\Phi_n(q)\mid q^n-1$, which divides $|S|\cdot
d$ ($d=\gcd(n,q-1)$); the Levi factors ($\GL_{n-1}(q)$-type and
$\GL_2{\times}\GL_{n-2}$-type) have cyclotomic exponents
$\leq(n-1)f<nf$; $x\mid2df$ and $r\equiv1\pmod{nf}$ rules out
$r\mid2f$ ($r>nf\geq3f>2f$) and $r\mid d$ (else $r\mid q-1$, so
$\ord_r(p)\mid f$). Zsigmondy exceptions $p^{nf}=2^6$:
$(n,f)=(3,2)$ is $\PSL(3,4)$ (machine base, all twelve $X$), and
$(n,f)=(6,1)$, $S=\PSL(6,2)$: take instead $r=31$ ($\ord_{31}(2)=5$)
and the pair $([P_2],[P_3])$, whose Levi types $\GL_2{\times}\GL_4$
and $\GL_3{\times}\GL_3$ have exponents $\leq4<5$, so
$31\nmid|P_2||P_3|$ while $v_{31}(|\PSL(6,2)|)=1$ and $x\mid2$.
\item $D_n(q)$, $n\geq4$: pair $([P_1],[P_2])$ (singular point and
singular line stabilizers), $r$ primitive for $p^{2(n-1)f}-1$:
$r\mid\Phi_{2n-2}(q)$, which divides
$q^{n(n-1)}(q^n-1)\prod_{i=1}^{n-1}(q^{2i}-1)$; Levi factors
($D_{n-1}$- and $A_1{\times}D_{n-2}$-type) have exponents
$\leq\max((n-1)f,2(n-2)f)<2(n-1)f$; $x\mid24f$ (diagonal $\leq4$,
field $f$, graph $\leq6$) and $r\equiv1\pmod{2(n-1)f}$ with
$2(n-1)\geq6$ gives $r>6f$ and $r\nmid24$ (as $r\equiv1\bmod6$ forces
$r\geq7$, $r\notin\{2,3\}$). Zsigmondy exception $p^{2(n-1)f}=2^6$:
$(n,f)=(4,1)$, $S=\Omega^+(8,2)$, handled for all $X$ in Case~B3.
\item $E_6(q)$: see Case~B4; the pair used there is stable for every
$X$.
\item $\Sp(4,2^f)$: pair $([P_1],[P_2])$, $r$ primitive for $2^{4f}-1$
(exists, $4f\geq8$): $r\mid\Phi_4(q)=q^2+1$; Levi factors $\Sp_2(q)$
and $\GL_2(q)$ have exponents $\leq2f<4f$; here $x\mid f$ and $r>4f$.
\item $F_4(2^f)$: pair $([P_1],[P_2])$, $r$ primitive for $2^{12f}-1$:
$r\mid\Phi_{12}(q)$; all proper Levis (types $B_3$, $C_3$,
$A_2{\times}A_1$, $A_1{\times}A_2$) have exponents $\leq6f<12f$;
$x\mid f<r$. (This completes the $p=2$ case excluded from
Theorem~\ref{thm:nograph}.)
\item $G_2(3^f)$: pair $([P_1],[P_2])$, $r$ primitive for $3^{6f}-1$:
$r\mid\Phi_6(q)$; Levis of type $A_1$ have exponents $\leq2f<6f$;
$x\mid f<r$. (The $p=3$ case of $G_2$.)
\end{itemize}

\emph{Case B: $X$ with nontrivial graph part.} Novelty classes below are
standard parabolics $Q_J$ handled by Lemma~\ref{lem:P}: each $[Q_J]$ is
$X$-stable because $J$ is invariant under the relevant diagram symmetry
(an automorphism with graph part $\gamma$ maps $[Q_J]$ to
$[Q_{\gamma(J)}]$), and its proper parabolic overgroups fall into
nontrivial $\gamma$-orbits, giving hypothesis (b) of Lemma~\ref{lem:P};
so $B_{Q_J}$ is maximal in $G$ by Lemma~\ref{lem:B}, and pairs of
distinct classes are non-conjugate by Lemma~\ref{lem:C}. The primes $r$
are those of Case~A unless stated; $r$-avoidance for $Q_J$ follows from
its own Levi order (a product of cyclotomic factors of the Levi of type
$J$ and powers of $q-1$), as listed.

\emph{B1. $\PSL(3,q)$, $q\geq5$} ($q\leq4$: machine base, including all
twelve $X$ for $\PSL(3,4)$). Pair: $[B]$ (novelty; the proper overgroups
$P_1,P_2$ are swapped by $\delta$, which exchanges points and lines of
$PG(2,q)$) and $[N]$, $N:=N_S(T)\cong((q-1)^2/d).S_3$ the normalizer of
a maximal split torus ($d=\gcd(3,q-1)$), maximal for $q\geq5$
\cite{Mitchell1911,Hartley1925}, \cite[Table~8.3]{BHRD} (class
$\mathcal{C}_2$), one class, preserved by every automorphism. Take
$r$ primitive for $p^{3f}-1$: $r\mid\Phi_3(q)=q^2+q+1$; $r\nmid|B|=
q^3(q-1)^2/d$ (primitivity) and $r\nmid|N|=6(q-1)^2/d$ (as
$r\equiv1\bmod3$ forces $r\geq7>6$, $r\nmid6$); $\vp{r}(|S|)\geq1$.
Also $x\mid2df$ and $r>3f$ rules out $r\mid2f$ and $r\mid d=3$.
Zsigmondy exception $p^{3f}=2^6$ is $\PSL(3,4)$: machine base.

\emph{B2. $\PSL(n,q)$, $n\geq4$.} For $n=4$: pair $([P_2],[Q_{\{2\}}])$.
Node $2$ is $\delta$-fixed, so $[P_2]$ is an $X$-stable \emph{maximal}
class (Levi $A_1{\times}A_1$, exponents $\leq2f$); $Q_{\{2\}}$ is the
stabilizer of an incident point--hyperplane flag, with proper overgroups
$P_1,P_3$ swapped by $\delta$ (Lemma~\ref{lem:P}); its Levi ($A_1$ plus
torus) has exponents $\leq2f$. Take $r$ primitive for $p^{4f}-1$: it
avoids both ($2f<4f$), $\vp{r}(|S|)\geq1$, and $x\mid2df$
($d=\gcd(4,q-1)$) with $r\equiv1\pmod{4f}$ excluding $r\mid2f$ and
$r\mid d$. No Zsigmondy exception ($2^{4f}=2^6$ impossible).
For $n\geq5$: pair of novelties $([Q_{J_1}],[Q_{J_2}])$,
$J_1:=\Delta\setminus\{1,n{-}1\}$, $J_2:=\Delta\setminus\{2,n{-}2\}$:
the stabilizers of an incident point--hyperplane flag and an incident
line--$(n{-}2)$-space flag. Both node sets $\{1,n{-}1\}$, $\{2,n{-}2\}$
are $\delta$-orbits, so both $[Q_{J_i}]$ are $X$-stable; the proper
overgroups of $Q_{J_1}$ are $P_1,P_{n-1}$ (swapped by $\delta$), of
$Q_{J_2}$ are $P_2,P_{n-2}$ (swapped, distinct classes since $n\geq5$).
Levi types: blocks $(1,n{-}2,1)$, i.e.\ $\GL_{n-2}$ plus torus,
exponents $\leq(n{-}2)f$; and blocks $(2,n{-}4,2)$, exponents
$\leq\max(2,n{-}4)f$. The prime $r$ of Case~A avoids both, and
$\vp{r}(x)=0$ as there. Zsigmondy exception $(n,f)=(6,1)$,
$S=\PSL(6,2)$: take $r=31$ again: the two Levi orders involve only
$\GL_4(2)$, $\GL_2(2)$ factors and $2$-powers, so $31$ divides neither,
while $v_{31}(|S|)=1>0=v_{31}(2)$.

\emph{B3. $D_n(q)$.} For $n\geq5$ the graph part is at most the $C_2$
swapping the two spinor nodes $n{-}1,n$ and fixing $1,\dots,n{-}2$: the
Case~A pair $([P_1],[P_2])$ is $X$-stable for \emph{every} $X$ and
Case~A's argument goes through unchanged (its $x$-bound already allowed
the graph factor). For $n=4$: if the graph part is trivial or the
normalized $C_2$ (fixing nodes $1,2$), the pair $([P_1],[P_2])$ is
$X$-stable and Case~A applies. If the graph part contains a triality,
take the pair $([P_2],[Q_{\{2\}}])$: node $2$ is fixed by all of $S_3$,
so $[P_2]$ is an $X$-stable maximal class (Levi
$A_1{\times}A_1{\times}A_1$, exponents $\leq2f$), and $Q_{\{2\}}$ (Levi
$A_1$ plus torus, exponents $\leq2f$) has proper overgroups
$Q_{\{1,2\}},Q_{\{2,3\}},Q_{\{2,4\}}$ and $P_4,P_3,P_1$, which fall
into the two regular $\langle\text{triality}\rangle$-orbits induced by
the permutation of the end nodes $\{1,3,4\}$; none is $X$-stable
(Lemma~\ref{lem:P}). The prime $r$ primitive for $p^{6f}-1$ avoids both
classes, $\vp{r}(|S|)\geq1$, and $r\nmid x$ as in Case~A.
\emph{Zsigmondy exception ($f=1$, $p=2$): $S=\Omega^+(8,2)$.} Here
$x\mid24$; take $r=5=\Phi_4(2)$, so
$v_5(|S|)=v_5\bigl(2^{12}(2^6-1)(2^4-1)^2(2^2-1)\bigr)=2$ and
$v_5(x)\leq v_5(24)=0$. Triality case, pair $([P_2],[Q_{\{2\}}])$: both
Levi orders are $\{2,3\}$-numbers, so $d=2-0-0=2>0$. Trivial/$C_2$
case, pair $([P_1],[P_2])$: $v_5(|P_1|)=v_5(|\SL_4(2)|)=1$ (Levi
$D_3\cong A_3$), $v_5(|P_2|)=0$, so $d=2-1-0=1>0$.

\emph{B4. $E_6(q)$, all $X$.} The diagram involution fixes nodes $2$
and $4$ (Bourbaki numbering: $\delta$ swaps $1\leftrightarrow6$,
$3\leftrightarrow5$). Pair $([P_2],[P_4])$: two $X$-stable classes of
\emph{maximal} parabolics for every $X$. $r$ primitive for $p^{12f}-1$:
$r\mid\Phi_{12}(q)$, which divides $|S|\cdot d$; the Levi of $P_2$ has
type $A_5$ (exponents $\leq6f$), that of $P_4$ type
$A_2{\times}A_1{\times}A_2$ ($\leq3f$), both $<12f$; $x\mid2df\leq6f<r$
($d=\gcd(3,q-1)$). No Zsigmondy exception. This simultaneously covers
Case~A for $E_6$.

\emph{B5. $\Sp(4,q)$, $q=2^f$, $f\geq2$, $X$ with graph part.} ($f=1$:
$\Sp(4,2)$ is not simple; its derived group $A_6$ is in the machine
base.) Pair: $[B]$ (novelty: the proper overgroups $P_1,P_2$ (the point
and totally-isotropic-line stabilizers) are swapped by $\rho$;
Lemma~\ref{lem:P}) and $[V]$, where
$V:=\Fix(\tilde\rho^{\,f})\cap S$ for $\tilde\rho$ the Steinberg
endomorphism of the ambient algebraic group with $\tilde\rho^2=
\varphi_2$ and $\tilde\rho^{\,2f}=F_q$. Explicitly $V=\Sz(q)$ for $f$
odd and $V=\Sp(4,2^{f/2})$ (subfield subgroup, index-two field
extension) for $f$ even; $V$ is maximal in both cases
\cite[Table~8.14]{BHRD}, \cite{Suzuki1962}. \emph{Stability of $[V]$:}
$\Aut(S)=\Inn(S)\langle\rho\rangle$ with $\rho=\tilde\rho|_S$;
$\tilde\rho$ commutes with its own power $\tilde\rho^{\,f}$, so
$\rho(V)=V$, and inner automorphisms preserve $[V]$ trivially; hence
$[V]$ is $\Aut$-stable. \emph{Primes:} $f$ odd, $2f\neq6$: $r$
primitive for $2^{2f}-1$, so $r\mid q+1$:
$\vp{r}(|S|)=\vp{r}\bigl((q^2-1)(q^4-1)\bigr)=2\vp{r}(q+1)\geq2$;
$\vp{r}(|B|)=\vp{r}(q^4(q-1)^2)=0$; $\vp{r}(|\Sz(q)|)=0$ since
$\gcd(q+1,\,q(q^2+1)(q-1))=1$ for $q$ even; $x\mid2f<r$. If $2f=6$
($q=8$): $r=3$ gives $v_3(|\Sp(4,8)|)=v_3(63\cdot4095)=4>1\geq
v_3(x)$, with $v_3(|B|)=v_3(|\Sz(8)|)=0$. $f$ even: $r$ primitive for
$2^{4f}-1$ (exists, $4f\geq8$), $r\mid q^2+1$: $\vp{r}(|S|)\geq1$;
$|\Sp(4,q_0)|=q_0^4(q_0^2-1)(q_0^4-1)$ with $q_0^2=q$ has cyclotomic
exponents $\leq2f<4f$, so $\vp{r}(|V|)=0$; $\vp{r}(|B|)=0$;
$x\mid2f<r$. ($\Sp(4,4)$ doubles as the machine template: the sweep-K3
certificate pair (Borel, $\Sp(4,2)$) at $p=5=\Phi_4(2)$ and
$p=17=\Phi_8(2)$ matches this construction exactly.)

\emph{B6. $F_4(q)$, $q=2^f$, $X$ with graph part.} $\rho$ reverses the
diagram $1$--$2$$\Rightarrow$$3$--$4$ (swaps $1\leftrightarrow4$,
$2\leftrightarrow3$). Pair of novelties $([Q_{J_1}],[Q_{J_2}])$,
$J_1:=\{2,3\}$, $J_2:=\{1,4\}$, both $\rho$-invariant node sets.
Proper overgroups of $Q_{J_1}$: $Q_{\{1,2,3\}}=P_4$ and
$Q_{\{2,3,4\}}=P_1$, swapped by $\rho$; of $Q_{J_2}$: $Q_{\{1,2,4\}}$
and $Q_{\{1,3,4\}}$, swapped by $\rho$ (distinct types, hence distinct
classes). Lemma~\ref{lem:P} applies to both. Levi types: $B_2$
($\cong\Sp_4(q)$, exponents $\leq4f$) and $A_1{\times}A_1$
($\leq2f$). $r$ primitive for $2^{12f}-1$ avoids both;
$\vp{r}(|S|)\geq1$; $x\mid2f<12f<r$. No exceptions.

\emph{B7. $G_2(q)$, $q=3^f$, $X$ with graph part.} $\rho$ swaps the two
nodes (long and short roots), $\rho^2=\varphi_3$. Pair: $[B]$ (novelty:
overgroups $P_1,P_2$ swapped by $\rho$; Lemma~\ref{lem:P}) and $[V]$,
$V:=\Fix(\tilde\rho^{\,f})\cap S$ exactly as in B5: $V={}^2G_2(q)$ for
$f$ odd, maximal in $G_2(q)$ \cite{Kleidman1988} (for $f=1$,
$V={}^2G_2(3)\cong P\Gamma L_2(8)$, maximal in $G_2(3)$ by the \textsc{Atlas}
and \cite{Kleidman1988}), and $V=G_2(3^{f/2})$ for $f$ even (subfield,
index-two extension, maximal \cite{Kleidman1988}). $[V]$ is
$\Aut$-stable by the commuting-endomorphism argument of B5
($\Out(G_2(3^f))=\langle\rho\rangle\cong C_{2f}$). \emph{Primes:} $f$
odd: $r$ primitive for $3^{3f}-1$ (exists for all $f\geq1$: $p=3\neq2$
excludes the $2^6$ exception; $f=1$ gives $r=13\mid3^3-1$). Then
$r\mid\Phi_3(q)=q^2+q+1$, which divides $q^6-1$, hence
$|S|=q^6(q^6-1)(q^2-1)$; $r\nmid|B|=q^6(q-1)^2$ (primitivity);
$r\nmid|{}^2G_2(q)|=q^3(q^3+1)(q-1)$ since $r\mid q^3-1$ and
$\gcd(q^3-1,q^3+1)=2\not\ni r$; $x\mid2f$ and $r\equiv1\pmod{3f}$ gives
$r\geq3f+1>2f$, so $\vp{r}(x)=0<1\leq\vp{r}(|S|)$. $f$ even: $r$
primitive for $3^{6f}-1$, so $r\mid\Phi_6(q)=q^2-q+1\mid q^3+1$:
$r\nmid|B|$; $r\nmid|G_2(3^{f/2})|=q^3(q^3-1)(q-1)$ (as $r\mid q^3+1$);
$x\mid2f<6f<r$. ($G_2(3)$ additionally carries a full machine
certificate for both $X$.)

In every case Theorem~\ref{thm:D} ($k\geq2$) and
Theorem~\ref{thm:Dprime} ($k=1$) exclude $S$.
\end{proof}

\section{Proof of the Main Theorem}\label{sec:proof}

\begin{proposition}[Coverage]\label{prop:coverage}
Every non-abelian finite simple group is excluded in the sense of
Convention~\ref{conv:excluded}.
\end{proposition}

\begin{proof}
By the classification of finite simple groups, $S$ is alternating,
of Lie type, or one of the $26$ sporadic groups; we run through the
list and name the theorem responsible.
\begin{itemize}
\item $A_n$ ($n\geq5$): Theorem~\ref{thm:an}.
\item $\PSL(2,q)$ ($q\geq4$): Theorem~\ref{thm:psl2}.
\item $\PSL(n,q)$ ($n\geq3$), $D_n(q)$ ($n\geq4$), $E_6(q)$:
Theorem~\ref{thm:graph}.
\item $\PSp(2n,q)$: $n\geq3$, or $n=2$ with $q$ odd:
Theorem~\ref{thm:nograph}; $\Sp(4,2^f)$, $f\geq2$:
Theorem~\ref{thm:graph} ($\Sp(4,2)$ is not simple).
\item $\Omega(2n{+}1,q)$, $q$ odd, $n\geq3$: Theorem~\ref{thm:nograph}
(for $q$ even, $B_n(q)\cong C_n(q)$, already listed;
$\Omega(5,q)\cong\PSp(4,q)$).
\item $E_7(q)$, $E_8(q)$; $F_4(q)$, $p\neq2$; $G_2(q)$, $p\neq3$:
Theorem~\ref{thm:nograph}. $F_4(2^f)$ and $G_2(3^f)$:
Theorem~\ref{thm:graph}. ($G_2(2)$ and ${}^2G_2(3)$ are not simple;
their derived groups $\PSU(3,3)$ and $\PSL(2,8)$ are covered.)
\item $\PSU(n,q)$: $n=3$: Theorem~\ref{thm:twisted1}; $n\geq4$:
Theorem~\ref{thm:twisted2}.
\item ${}^2D_n(q)$ ($n\geq4$), ${}^3D_4(q)$, ${}^2E_6(q)$,
${}^2F_4(2^f)'$: the first three and ${}^2F_4(2^f)$ ($f$ odd $\geq3$)
by Theorem~\ref{thm:twisted2}; the Tits group ${}^2F_4(2)'$ by
Proposition~\ref{prop:sporadic}. (${}^2D_2$, ${}^2D_3$ are
$\PSL(2,q^2)$, $\PSU(4,q)$, already listed.)
\item $\Sz(2^f)$ ($f$ odd $\geq3$), ${}^2G_2(3^f)$ ($f$ odd $\geq3$):
Theorem~\ref{thm:twisted1}.
\item Sporadic groups: Proposition~\ref{prop:sporadic}.
\end{itemize}
Every finite simple group appears (repetitions from exceptional
isomorphisms are harmless), so the proposition follows.
\end{proof}

\begin{proof}[Proof of Theorem~\ref{thm:main}]
Suppose not, and let $G$ be a non-soluble group with property
$\mathrm{P}$ of least order. By Proposition~\ref{prop:min}, $G$ has a
unique minimal normal subgroup $N=S^k$, $S$ non-abelian simple, and
after conjugation $N\leq G\leq X\wr S_k$ as in Convention~\ref{conv:X}
(for $k=1$ this reads $S\leq G\leq\Aut(S)$ with $X=G$). By
Proposition~\ref{prop:coverage} there are classes
$[V]\neq[W]$, maximal in $\Pos_X(S)$ and normally saturating, and a
prime $p$ with $\vp{p}(|S|)-\vp{p}(|V|)-\vp{p}(|W|)>\vp{p}(x)$. By
Theorem~\ref{thm:D} (if $k\geq2$) or Theorem~\ref{thm:Dprime} (if
$k=1$), $G$ does not have property $\mathrm{P}$, a contradiction.
\end{proof}

\begin{corollary}
No almost simple group has property $\mathrm{P}$
(Tikhonenko--Tyutyanov \cite{TT2010}).
\end{corollary}

\begin{proof}
Theorem~\ref{thm:Dprime} with Proposition~\ref{prop:coverage}.
\end{proof}

\begin{remark}
The proof yields slightly more than solubility of the minimal
counterexample: for \emph{every} finite group $G$ with a unique minimal
normal subgroup $S^k$ ($S$ non-abelian simple), some pair of
non-conjugate maximal subgroups of $G$ fails to factorize $G$. The
reduction of \S\ref{sec:reduction} is needed only to pass from an
arbitrary non-soluble group to this configuration.
\end{remark}

\subsection*{Data Availability}
All code, GAP scripts, and log certificates required to reproduce the
machine-verified claims of \S\ref{sec:machine}, together with the
arithmetic receipts for the family proofs of \S\ref{sec:families}, are
hosted at
\texttt{https://github.com/RichieSater/kourovka-10-34} \cite{Repo}; its
\texttt{README} documents how to re-run every sweep and receipt.

\begin{thebibliography}{99}

\bibitem{BHRD}
J.~N.~Bray, D.~F.~Holt, C.~M.~Roney-Dougal,
\emph{The Maximal Subgroups of the Low-Dimensional Finite Classical
Groups}, London Math.\ Soc.\ Lecture Note Ser.\ 407, Cambridge Univ.\
Press, 2013.

\bibitem{Carter}
R.~W.~Carter, \emph{Simple Groups of Lie Type}, Wiley, London, 1972.

\bibitem{ATLAS}
J.~H.~Conway, R.~T.~Curtis, S.~P.~Norton, R.~A.~Parker, R.~A.~Wilson,
\emph{Atlas of Finite Groups}, Oxford Univ.\ Press, 1985.

\bibitem{CTblLib}
T.~Breuer, \emph{The GAP Character Table Library}, GAP package,
\texttt{https://www.math.rwth-aachen.de/\~{}Thomas.Breuer/ctbllib}.

\bibitem{Dickson}
L.~E.~Dickson, \emph{Linear Groups with an Exposition of the Galois
Field Theory}, Teubner, Leipzig, 1901.

\bibitem{GAP}
The GAP Group, \emph{GAP -- Groups, Algorithms, and Programming},
Version 4.16.0, 2026, \texttt{https://www.gap-system.org}.

\bibitem{GLS3}
D.~Gorenstein, R.~Lyons, R.~Solomon, \emph{The Classification of the
Finite Simple Groups, Number 3}, Math.\ Surveys Monogr.\ 40.3, Amer.\
Math.\ Soc., 1998.

\bibitem{Hartley1925}
R.~W.~Hartley, Determination of the ternary collineation groups whose
coefficients lie in the $GF(2^n)$, \emph{Ann.\ of Math.} \textbf{27}
(1925), 140--158.

\bibitem{Huppert}
B.~Huppert, \emph{Endliche Gruppen I}, Grundlehren Math.\ Wiss.\ 134,
Springer, Berlin, 1967.

\bibitem{Kleidman1988}
P.~B.~Kleidman, The maximal subgroups of the Chevalley groups $G_2(q)$
with $q$ odd, the Ree groups ${}^2G_2(q)$, and their automorphism
groups, \emph{J.\ Algebra} \textbf{117} (1988), 30--71.

\bibitem{Kourovka}
E.~I.~Khukhro, V.~D.~Mazurov (eds.), \emph{The Kourovka Notebook:
Unsolved Problems in Group Theory}, 21st ed., Sobolev Inst.\ Math.,
Novosibirsk, 2026; Problem 10.34 (V.~S.~Monakhov).

\bibitem{LPS1987}
M.~W.~Liebeck, C.~E.~Praeger, J.~Saxl, A classification of the maximal
subgroups of the finite alternating and symmetric groups,
\emph{J.\ Algebra} \textbf{111} (1987), 365--383.

\bibitem{Mitchell1911}
H.~H.~Mitchell, Determination of the ordinary and modular ternary
linear groups, \emph{Trans.\ Amer.\ Math.\ Soc.} \textbf{12} (1911),
207--242.

\bibitem{Nagura}
J.~Nagura, On the interval containing at least one prime number,
\emph{Proc.\ Japan Acad.} \textbf{28} (1952), 177--181.

\bibitem{Repo}
R.~Sater, \emph{Kourovka 10.34: proofs, GAP scripts and machine
certificates}, GitHub repository,
\texttt{https://github.com/RichieSater/kourovka-10-34}, 2026.

\bibitem{Steinberg}
R.~Steinberg, Automorphisms of finite linear groups,
\emph{Canad.\ J.\ Math.} \textbf{12} (1960), 606--615.

\bibitem{Suzuki1962}
M.~Suzuki, On a class of doubly transitive groups,
\emph{Ann.\ of Math.} \textbf{75} (1962), 105--145.

\bibitem{TT2010}
T.~V.~Tikhonenko, V.~N.~Tyutyanov, On factorization of finite almost
simple groups by nonconjugate maximal subgroups, \emph{Sib.\ Math.\ J.}
\textbf{51}, no.~1 (2010), 174--177; DOI 10.1007/s11202-010-0018-3.

\bibitem{Zsigmondy}
K.~Zsigmondy, Zur Theorie der Potenzreste, \emph{Monatsh.\ Math.\
Phys.} \textbf{3} (1892), 265--284.

\end{thebibliography}

\end{document}
